% ============================================================
% PAPER A -- FOUNDATIONS APPENDIX
% "Foundations of the Canonical MASA-Stack Class"
% A rigorous, self-contained proof appendix for the
% geometric/cohomological layer of Paper A.
%
% Every result is labelled explicitly as one of:
%   [COMPLETE]  a complete proof, checkable by a specialist
%               in equivariant topology / stack cohomology
%               without any repository script;
%   [SKETCH]    a proof reduced to a named, precisely cited
%               standard theorem whose hypotheses are verified;
%   [GAP]       an honestly stated open sub-lemma.
% A "Status of each step" table closes the document.
%
% Self-contained: manual thebibliography, no BibTeX.
% ============================================================
\documentclass[11pt]{article}
\usepackage[margin=1in]{geometry}
\usepackage{amsmath,amssymb,amsthm,mathtools}
\usepackage{booktabs}
\usepackage[hidelinks]{hyperref}
\usepackage{lmodern}

\newtheorem{theorem}{Theorem}[section]
\newtheorem{proposition}[theorem]{Proposition}
\newtheorem{lemma}[theorem]{Lemma}
\newtheorem{corollary}[theorem]{Corollary}
\theoremstyle{definition}
\newtheorem{definition}[theorem]{Definition}
\newtheorem{remark}[theorem]{Remark}
\newtheorem{convention}[theorem]{Convention}

\newcommand{\U}{\mathrm{U}}
\newcommand{\PU}{\mathrm{PU}}
\newcommand{\PN}{\mathrm{PN}}
\newcommand{\SU}{\mathrm{SU}}
\newcommand{\Nt}{N(T)}
\newcommand{\Z}{\mathbb{Z}}
\newcommand{\R}{\mathbb{R}}
\newcommand{\C}{\mathbb{C}}
\newcommand{\HH}{H^{2}}
\newcommand{\Hom}{\mathrm{Hom}}
\newcommand{\Ext}{\mathrm{Ext}}
\newcommand{\im}{\mathrm{im}}
\newcommand{\res}{\mathrm{res}}
\newcommand{\infl}{\mathrm{inf}}
\newcommand{\id}{\mathrm{id}}
\newcommand{\Wn}{\overline{W}_{n}}
\newcommand{\That}{\widehat{T}}
\newcommand{\Tp}{T'}
\newcommand{\Thatp}{\widehat{T'}}
\newcommand{\gp}{\mathrm{gp}}
\newcommand{\sm}{\mathrm{sm}}
\newcommand{\ltimesop}{\ltimes}
\DeclareMathOperator{\Aut}{Aut}

% status tags
\newcommand{\STAT}[1]{\noindent\textbf{Status: #1}}
\newcommand{\COMPLETE}{\textsc{Complete proof}}
\newcommand{\SKETCH}{\textsc{Proof sketch (named standard result)}}
\newcommand{\GAP}{\textsc{Gap (open sub-lemma)}}

\title{Foundations of the Canonical MASA--Stack Class\\
\large A rigorous proof appendix for the geometric/cohomological layer of Paper~A}
\author{Manuel Flores Gordillo\thanks{ORCID \texttt{0009-0006-8246-0343};
\texttt{manuel.flores@columbia.edu}. This appendix is self-contained and is written so
that the steps marked \textsc{Complete} can be verified without consulting any repository
script. Companion: \emph{The Canonical Class of the MASA Context Stack and the
Peres--Mermin Obstruction} (Paper~A).}}
\date{July 2026}

\begin{document}
\maketitle

\begin{abstract}
This appendix supplies rigorous foundations for the geometric and cohomological
constructions of Paper~A. We (A) fix the objects --- the space $X$ of maximal abelian
subalgebras (MASAs) of $M_n(\C)$ as the homogeneous space $\PU(n)/\PN(T)$, the action
groupoid $\PU(n)\ltimes X$ presenting the stack $[X/\PU(n)]$, and a single, explicitly
named cohomology theory that is never silently switched; (B) construct the canonical class
from $1\to\U(1)\to\U(n)\to\PU(n)\to1$ and prove its vacuity on $\U(n)\ltimes X$ and its
nontriviality on $\PU(n)\ltimes X$; (C) prove the order is exactly $n$, giving both
$n[\alpha]=0$ (determinant-twist splitting) and $k[\alpha]\neq0$ for $1\le k<n$, with the
role of the $S_n$-equivariant homomorphism lemma made explicit and the ``exotic splitting''
objections addressed; (D) state and prove the Morita-equivalence step, naming the theorem
and verifying its hypotheses; and (E) compute $H^2(\PN(T);\U(1))\cong \Z/n\oplus\Z_2$
($n\ge4$) via a Lyndon--Hochschild--Serre spectral sequence with all differentials,
extension problems, coefficient conventions, topology assumptions, and low-$n$ exceptions
recorded. \textbf{Intellectual honesty is the point:} every result is tagged
\textsc{Complete}, \textsc{Sketch} (reduced to a cited theorem whose hypotheses we verify),
or \textsc{Gap} (an explicit open sub-lemma), and a status table closes the document.
\end{abstract}

\tableofcontents

\section*{Reading guide and honesty protocol}
\addcontentsline{toc}{section}{Reading guide and honesty protocol}
Each numbered result carries a \textbf{Status} line immediately after its proof:
\begin{itemize}
\item \COMPLETE{}: a full argument, checkable by a specialist in equivariant topology /
stack cohomology using only standard textbook facts (cited), with no appeal to any
repository computation.
\item \SKETCH{}: the result is reduced to a precisely named theorem in the literature; we
verify that theorem's hypotheses but do not reprove it.
\item \GAP{}: an honestly flagged open sub-lemma; we state exactly what must be proved.
\end{itemize}
Two analytic inputs are used repeatedly and cited once and for all: the
Segal--Mitchison/Moore theory of the cohomology of topological groups
(\cite{Segal,MooreIII,MooreIV,Wigner}) and the Behrend--Xu / Tu--Xu--Laurent-Gengoux theory
of $S^1$-central extensions of Lie groupoids and its Morita invariance
(\cite{BehrendXu,TXL,Crainic}). Where a step depends on one of these it is marked
\textsc{Sketch} and the precise statement invoked is quoted.

% ======================================================================
\section{Item A: objects and the cohomology theory}\label{sec:A}
% ======================================================================

\subsection{The space of MASAs as a homogeneous space}

Fix $n\ge2$ and $H=\C^n$ with its standard Hermitian inner product. Write
$M_n(\C)=B(H)$.

\begin{definition}[MASA]
A \emph{maximal abelian subalgebra} (MASA) of $M_n(\C)$ is a unital $*$-subalgebra
$A\subseteq M_n(\C)$ that is commutative and is not properly contained in any commutative
unital $*$-subalgebra. Let $X$ denote the set of all MASAs of $M_n(\C)$.
\end{definition}

\begin{lemma}[Structure and homogeneity of $X$]\label{lem:homog}
Every MASA $A\subseteq M_n(\C)$ equals the algebra of operators diagonal in some
orthonormal basis; equivalently $A$ is determined by the associated orthogonal
decomposition $H=\bigoplus_{i=1}^n \ell_i$ into $n$ pairwise orthogonal lines (its minimal
projections). Consequently:
\begin{enumerate}
\item[(i)] The conjugation action $U\cdot A := UAU^{\dagger}$ of $\U(n)$ on $X$ is
transitive.
\item[(ii)] Let $D=\{\,\mathrm{diag}(d_1,\dots,d_n)\,\}$ be the standard diagonal MASA. Its
stabilizer in $\U(n)$ is the group $\Nt$ of \emph{monomial} unitaries (one nonzero entry
per row and column), i.e. $\Nt=T\rtimes S_n$ where $T=\U(1)^n$ is the diagonal maximal
torus and $S_n$ acts by permutation matrices. This is the normalizer of $T$ in $\U(n)$.
\item[(iii)] The centre $\U(1)=\{\lambda I\}$ acts trivially, so the action factors through
$\PU(n)=\U(n)/\U(1)$, and
\[
X\;\cong\;\U(n)/\Nt\;\cong\;\PU(n)/\PN(T),\qquad \PN(T):=\Nt/\U(1).
\]
In particular $X$ is a compact connected smooth manifold of dimension
$\dim\PU(n)-\dim\PN(T)=(n^2-1)-(n-1)=n^2-n$.
\end{enumerate}
\end{lemma}

\begin{proof}
A commutative $*$-subalgebra of $M_n(\C)$ is simultaneously diagonalizable by the spectral
theorem; maximality forces it to contain a complete system of $n$ rank-one orthogonal
projections and hence to be exactly the diagonal algebra in the corresponding orthonormal
basis. Two such algebras are conjugate by the unitary carrying one orthonormal basis to the
other, giving (i). The stabilizer of $D$ consists of unitaries permuting the coordinate
lines $\{\C e_i\}$ up to phase, i.e. monomial unitaries; a monomial unitary is uniquely
$P_\sigma\,\mathrm{diag}(t_1,\dots,t_n)$ with $\sigma\in S_n$, $t_i\in\U(1)$, giving
$\Nt=T\rtimes S_n$ and the classical identification $\Nt=N_{\U(n)}(T)$, proving (ii). The
scalars $\lambda I$ fix every subspace, hence act trivially; the induced $\PU(n)$-action is
still transitive with stabilizer $\Nt/\U(1)=\PN(T)$, giving (iii). Smoothness and the
dimension count are standard for coset spaces of compact Lie groups.
\end{proof}
\STAT{\COMPLETE.} (Spectral theorem; elementary Lie theory of $N_{\U(n)}(T)$.)

\begin{convention}[The torus quotient]\label{conv:torus}
Throughout, $\Tp:=T/\U(1)$ where $\U(1)\hookrightarrow T=\U(1)^n$ is the diagonal
$\lambda\mapsto(\lambda,\dots,\lambda)$. Then $\Tp\cong\U(1)^{n-1}$ is a torus of rank
$n-1$, the diagonal scalars being $S_n$-invariant we have a split extension
\begin{equation}\label{eq:PNsplit}
1\longrightarrow \Tp \longrightarrow \PN(T)=\Tp\rtimes S_n \longrightarrow S_n
\longrightarrow 1 ,
\end{equation}
the splitting $S_n\hookrightarrow\PN(T)$ being the images of the permutation matrices
$P_\sigma$. The (Pontryagin) character lattice of $\Tp$ is
\[
\Thatp=\Hom_{\mathrm{cont}}(\Tp,\U(1))=\{a\in\Z^n:\textstyle\sum_i a_i=0\}=A_{n-1},
\]
the root lattice of type $A_{n-1}$, with $S_n$ acting by permuting coordinates.
\end{convention}

\subsection{The action groupoid, the stack, and the chosen cohomology theory}

\begin{definition}[Action groupoid and stack]
For a Lie group $G$ acting smoothly on a manifold $M$, the \emph{action groupoid}
$G\ltimes M$ has object manifold $M_0=M$ and arrow manifold $M_1=G\times M$, with
$(g,x)\colon x\to g\cdot x$, composition $(h,g\cdot x)(g,x)=(hg,x)$, units $(e,x)$, and
inverses $(g,x)^{-1}=(g^{-1},g\cdot x)$. It presents the quotient differentiable stack
$[M/G]$. We write
\[
G:=\U(n)\ltimes X,\qquad \bar G:=\PU(n)\ltimes X ,
\]
so that $[X/\PU(n)]$ is the \emph{MASA context stack}.
\end{definition}

We must now fix, once and for all, a single cohomology theory. The requirements are: (a)
degree $2$ must classify \emph{central extensions by $\U(1)$} (of the group, resp.\ of the
groupoid), so that ``the class of an extension'' is literally an $H^2$; (b) it must be
Morita invariant for Lie groupoids (Item~D); (c) on a Lie group it must reduce to a theory
in which the computations of Items C and E are the standard ones.

\begin{definition}[The cohomology theory used, and its coefficients]\label{def:coh}
Let $\underline{\U(1)}$ denote the sheaf of smooth $\U(1)$-valued functions, regarded as a
groupoid module with \emph{trivial} action (all our extensions are central). For a Lie
groupoid $\mathcal G=(\mathcal G_1\rightrightarrows \mathcal G_0)$ we use the
\emph{Segal--Mitchison/Behrend--Xu groupoid cohomology}
$H^{\bullet}(\mathcal G;\underline{\U(1)})$, defined as the cohomology of the double complex
obtained from the nerve $N_{\bullet}\mathcal G$ (the simplicial manifold with
$N_p\mathcal G=\mathcal G_1\times_{\mathcal G_0}\cdots\times_{\mathcal G_0}\mathcal G_1$)
and a Segal--Mitchison soft resolution of $\underline{\U(1)}$; see
\cite[\S2--3]{BehrendXu} and \cite[\S2]{TXL}. Coefficients are written multiplicatively.
We record three properties, all standard in the cited references:
\begin{enumerate}
\item[(C1)] \textbf{(Extensions.)} $H^2(\mathcal G;\underline{\U(1)})$ is in natural
bijection with the group of isomorphism classes of \emph{$S^1$-central extensions}
(equivalently, $\U(1)$-gerbes with connective structure discarded) of $\mathcal G$; Baer
sum corresponds to addition \cite[Thm.~3.9]{BehrendXu}.
\item[(C2)] \textbf{(Group case.)} If $\mathcal G=G$ is a Lie group viewed as a one-object
groupoid, then $H^{\bullet}(G;\underline{\U(1)})$ equals the Segal--Mitchison smooth group
cohomology $H^{\bullet}_{\sm}(G;\U(1))$ \cite{Segal,MooreIII}. If moreover $G$ is finite
(discrete), it equals ordinary group cohomology $H^{\bullet}(G;\U(1))$ with $\U(1)$ a
trivial module.
\item[(C3)] \textbf{(Morita invariance.)} A Morita equivalence of Lie groupoids induces an
isomorphism on $H^{\bullet}(-;\underline{\U(1)})$ \cite[Thm.~3.9, Prop.~3.10]{BehrendXu},
\cite[\S2]{TXL}, \cite[Thm.~4]{Crainic}.
\end{enumerate}
Whenever ``$H^2(-;\U(1))$'' appears below it means \emph{this} theory (equivalently, on
Lie groups, $H^2_{\sm}$; on finite groups, ordinary $H^2$). We never switch silently.
\end{definition}
\STAT{\COMPLETE{} as a definition; the three properties (C1)--(C3) are \SKETCH{} (they are
the cited theorems of Behrend--Xu, Tu--Xu--Laurent-Gengoux, Segal--Mitchison, Moore, and
we invoke rather than reprove them).}

\begin{remark}[Why this theory and not another]
The alternatives are equivalent for our inputs but each obscures one requirement:
Borel/Moore measurable cohomology $H^{\bullet}_{\mathrm{meas}}(G;\U(1))$ agrees with
$H^{\bullet}_{\sm}$ for the (compact, or Lie) groups here \cite{MooreIII,WagnerAustin} but
has no built-in stack meaning; Segal--Borel equivariant cohomology $H^{\bullet}_{G}(X;\Z)$
computes the \emph{topological} invariants of $[X/G]$ but sits in degree shifted by one
relative to $\U(1)$-central extensions. Definition~\ref{def:coh} is the unique choice for
which (a) the ``class of a central extension'' is directly $H^2$, and (b) Morita invariance
(Item~D) is a cited theorem. All numerical outputs below are theory-independent because
they are computed after Morita reduction to the compact Lie group $\PN(T)$, where
$H^{\bullet}_{\sm}=H^{\bullet}_{\mathrm{meas}}$.
\end{remark}

% ======================================================================
\section{Item B: the canonical class, its vacuity, and its nontriviality}\label{sec:B}
% ======================================================================

\subsection{Construction from the central extension of $\PU(n)$}

The central extension
\begin{equation}\label{eq:UPU}
1\longrightarrow \U(1)\longrightarrow \U(n)\xrightarrow{\ \pi\ }\PU(n)\longrightarrow 1
\end{equation}
has a class $[\pi]\in H^2_{\sm}(\PU(n);\U(1))$ (property (C1) for the one-object groupoid
$\PU(n)$). Pick a Borel section $s\colon\PU(n)\to\U(n)$ of $\pi$; the associated normalized
$2$-cocycle is
\[
c_s(\bar V,\bar U)\;=\;s(\bar V)\,s(\bar U)\,s(\overline{VU})^{-1}\ \in\ \U(1),
\qquad \bar V,\bar U\in\PU(n).
\]

\begin{definition}[The canonical class]\label{def:canon}
Let $F\colon \bar G=\PU(n)\ltimes X\to \PU(n)$ be the smooth functor sending the arrow
$(\,\bar U, C\,)\colon C\to \bar U C\bar U^{\dagger}$ to $\bar U$ (and the unique object to
the unique object). The \emph{canonical class} of the MASA context stack is
\[
[\omega]\;:=\;F^{*}[\pi]\ \in\ H^2(\bar G;\underline{\U(1)}),
\]
represented by the groupoid $2$-cocycle
$\omega\big((\bar V,\cdot),(\bar U,\cdot)\big)=c_s(\bar V,\bar U)$, constant along $X$.
Similarly define $F_0\colon G=\U(n)\ltimes X\to\U(n)$, and note $F=\bar\pi\circ F_0$ where
$\bar\pi\colon\U(n)\to\PU(n)$ on labels.
\end{definition}

\subsection{Vacuity on $\U(n)\ltimes X$}

\begin{theorem}[Vacuity]\label{thm:vacuity}
$F_0^{*}[\pi]=[\omega]|_{G}=0$ in $H^2(G;\underline{\U(1)})$, $G=\U(n)\ltimes X$.
\end{theorem}

\begin{proof}
On $G$ every arrow carries an honest unitary $U\in\U(n)$ (not merely its class). Define a
$1$-cochain $\lambda\colon G_1=\U(n)\times X\to\U(1)$ by
\[
\lambda(U,C)\;:=\;U\,s(\bar U)^{-1}\in\U(1),
\]
which lies in $\U(1)$ because $U$ and $s(\bar U)$ have the same image $\bar U$ in $\PU(n)$,
so their ratio is central. The pulled-back cocycle is
$\omega|_G\big((V,\cdot),(U,\cdot)\big)=s(\bar V)s(\bar U)s(\overline{VU})^{-1}$. Compute
the groupoid coboundary of $\lambda$ (trivial coefficient action):
\[
(\delta\lambda)\big((V,\cdot),(U,\cdot)\big)
=\lambda(V,\cdot)\,\lambda(U,\cdot)\,\lambda(VU,\cdot)^{-1}
= \big(Vs(\bar V)^{-1}\big)\big(Us(\bar U)^{-1}\big)\big(VU\,s(\overline{VU})^{-1}\big)^{-1}.
\]
Since the factors are scalars they commute, and $V,U$ cancel against $(VU)^{-1}$:
\[
(\delta\lambda)=s(\bar V)^{-1}s(\bar U)^{-1}s(\overline{VU})
=\big(s(\bar V)s(\bar U)s(\overline{VU})^{-1}\big)^{-1}=\omega|_G^{-1}.
\]
Hence $\omega|_G=\delta(\lambda^{-1})$ is a coboundary and $[\omega]|_G=0$. Equivalently
$(U,C)\mapsto\big(\lambda(U,C),(U,C)\big)$ is a strict homomorphic section of the pulled-back
central extension groupoid.
\end{proof}
\STAT{\COMPLETE.} The mechanism is precisely that $F_0$ factors the classifying functor
through $\U(n)$, along which \eqref{eq:UPU} is (tautologically) split by the identity; the
extra arrow data of the groupoid supplies the honest lift $U$.

\begin{remark}[The cocycle is not itself symmetric]
The vanishing is a statement about the \emph{groupoid} class, not about $c_s$ as a function
on $\PU(n)\times\PU(n)$: for generic $s$ the cocycle $c_s$ disagrees with its transpose and
is not a group coboundary. Theorem~\ref{thm:vacuity} says only that on $\U(n)\ltimes X$ the
class dies; the construction must therefore be run on $\bar G=\PU(n)\ltimes X$, over which
the centre acts trivially and $F$ does not factor through $\U(n)$.
\end{remark}

\subsection{Nontriviality on $\PU(n)\ltimes X$}

\begin{theorem}[Nontriviality]\label{thm:nontrivial}
For every $n\ge2$, $[\omega]\neq0$ in $H^2(\bar G;\underline{\U(1)})$; indeed $[\omega]$ has
order exactly $n$.
\end{theorem}

\begin{proof}
By the Morita reduction (Proposition~\ref{prop:morita}, Item~D) there is an isomorphism
$H^{\bullet}(\bar G;\underline{\U(1)})\cong H^{\bullet}(\PN(T);\U(1))$ carrying $[\omega]$
to the class $[\alpha]$ of the extension \eqref{eq:ext1}. By the order theorem
(Theorem~\ref{thm:order}, Item~C) $[\alpha]$ has order exactly $n\ge2$, hence is nonzero,
hence so is $[\omega]$.
\end{proof}
\STAT{\COMPLETE{} \emph{given} Items C and D} (on which it depends and which are proved
below without circularity: Item~D is a Morita statement independent of the order, Item~C is
a finite-group computation).

% ======================================================================
\section{Item D: the Morita reduction, with hypotheses verified}\label{sec:D}
% ======================================================================

\begin{theorem}[Named theorem: Morita equivalence of a transitive action groupoid with a
stabilizer]\label{thm:MoritaGeneral}
Let a Lie group $G$ act smoothly and transitively on a manifold $M$, and let $H=G_{x_0}$ be
the (closed) stabilizer of a point $x_0\in M$, so $M\cong G/H$. Then the action groupoid
$G\ltimes M$ is Morita equivalent to $H$ (a one-object groupoid), with principal bibundle
$G$ itself. Consequently $H^{\bullet}(G\ltimes M;\underline{\U(1)})\cong
H^{\bullet}(H;\underline{\U(1)})$.
\end{theorem}

This is standard \cite[Ch.~5]{MoerdijkMrcun}, \cite[\S2.4]{delHoyo}; the cohomological
consequence is (C3). We must verify its hypotheses \emph{for our data} and identify the
class.

\begin{proposition}[Morita reduction of the canonical class]\label{prop:morita}
The bibundle $P:=\PU(n)$ realises a Morita equivalence
$\bar G=\PU(n)\ltimes X\ \simeq_{\mathrm{Mor}}\ \PN(T)$, and under the induced isomorphism
$H^{\bullet}(\bar G;\underline{\U(1)})\cong H^{\bullet}(\PN(T);\U(1))$ the canonical class
$[\omega]$ corresponds to the class $[\alpha]$ of the central extension
\begin{equation}\label{eq:ext1}
1\longrightarrow \U(1)\longrightarrow \Nt \longrightarrow \PN(T)\longrightarrow 1 .
\end{equation}
\end{proposition}

\begin{proof}
\emph{Hypotheses of Theorem~\ref{thm:MoritaGeneral}.} By Lemma~\ref{lem:homog} the
$\PU(n)$-action on $X$ is smooth and transitive with closed stabilizer $\PN(T)$ at the base
MASA $D$, and $X\cong\PU(n)/\PN(T)$. All groups are compact Lie, so the stabilizer is
closed and the quotient map is a smooth principal $\PN(T)$-bundle. Hence
Theorem~\ref{thm:MoritaGeneral} applies with $G=\PU(n)$, $M=X$, $H=\PN(T)$.

\emph{The bibundle, explicitly.} Let $P=\PU(n)$ with:
(i) a right $\PN(T)$-action by right translation $p\cdot h=ph$, which is free and proper
(right translation in a compact Lie group by a closed subgroup) with quotient
$P/\PN(T)=X$ and moment map $\tau\colon P\to X$, $\tau(p)=p\cdot D:=pDp^{\dagger}$;
(ii) a left $\bar G$-action over $\tau$: an arrow $(\bar U,x)$ with $x=\tau(p)$ acts by
$(\bar U,x)\cdot p=\bar U p$; this is well defined since $\tau(\bar U p)=\bar U\cdot\tau(p)$,
and it is \emph{principal}, i.e.\ $\tau$ is a surjective submersion and $\bar G$ acts freely
and transitively on the $\tau$-fibres (given $p,p'$ with $\tau(p)=\tau(p')$ the unique arrow
is $(\,p'p^{-1},\tau(p)\,)$). The two actions commute:
$(\bar U p)h=\bar U(ph)$. Thus $P$ is a $\bar G$--$\PN(T)$ principal bibundle, which is the
definition of a Morita equivalence \cite[\S5.4]{MoerdijkMrcun}. Property (C3) then gives the
cohomology isomorphism.

\emph{Transport of the class.} The classifying functor $F\colon\bar G\to\PU(n)$ of
Definition~\ref{def:canon} restricts, under the equivalence
$\bar G\simeq_{\mathrm{Mor}}\PN(T)$, to the inclusion of isotropy
$\iota\colon\PN(T)\hookrightarrow\PU(n)$ (an arrow of $\bar G$ fixing $D$ has label in the
stabilizer $\PN(T)$, and $F$ reads off that label). Naturality of pullback gives
$[\omega]=F^{*}[\pi]\ \mapsto\ \iota^{*}[\pi]$ under the isomorphism. Finally $\iota^{*}[\pi]$
is the class of the pullback of \eqref{eq:UPU} along
$\iota\colon\PN(T)\hookrightarrow\PU(n)$, namely of
\[
1\to \U(1)\to \pi^{-1}(\PN(T))\to \PN(T)\to1 .
\]
Because $\U(1)=\ker\pi\subseteq \Nt$ (scalars are monomial) and $\Nt/\U(1)=\PN(T)$, we have
$\pi^{-1}(\PN(T))=\Nt$, so the pullback is exactly \eqref{eq:ext1}. Hence
$[\omega]\leftrightarrow[\alpha]$.
\end{proof}
\STAT{\COMPLETE{} for the bibundle construction and the transport-of-class identification;
\SKETCH{} for the single invocation of (C3) (the Morita invariance theorem
\cite{BehrendXu,TXL,Crainic}), whose hypotheses (both actions principal, actions commute)
are verified above.} No step here uses ``by Morita invariance'' without exhibiting the
principal bibundle and checking its axioms.

\begin{remark}[On the naturality step]
The only non-mechanical point is that $F$ corresponds to $\iota$ under the equivalence. This
is because a Morita equivalence $\bar G\simeq H$ identifies the category of $\bar
G$-representations/cocycles with that of $H$, and the isotropy inclusion $H\hookrightarrow
\bar G$ (as the automorphisms of the base object) composed with $F$ is $\iota$. For the
$H^2$-level statement (C1) this is the assertion that the $S^1$-central extension classified
by $[\omega]$ restricts on isotropy to \eqref{eq:ext1}; see \cite[Prop.~3.10]{BehrendXu}.
\end{remark}

% ======================================================================
\section{Item C: the order is exactly $n$}\label{sec:C}
% ======================================================================

By Proposition~\ref{prop:morita} it suffices to compute the order of $[\alpha]\in
H^2(\PN(T);\U(1))$, the class of \eqref{eq:ext1}. We prove $n[\alpha]=0$ and
$k[\alpha]\neq0$ for $1\le k<n$.

\subsection{Upper bound: $n[\alpha]=0$ via the determinant twist}

\begin{lemma}[Determinant on $\Nt$]\label{lem:det}
$\det\colon\Nt\to\U(1)$, $\det(P_\sigma\,\mathrm{diag}(t_1,\dots,t_n))=\mathrm{sgn}(\sigma)
\prod_i t_i$, is a continuous homomorphism with $\det(\lambda g)=\lambda^{n}\det(g)$ for
$\lambda\in\U(1)$, $g\in\Nt$.
\end{lemma}
\begin{proof}
The determinant is multiplicative and continuous on $\U(n)\supseteq\Nt$; homogeneity is
$\det(\lambda I\cdot g)=\lambda^n\det g$.
\end{proof}
\STAT{\COMPLETE.}

\begin{theorem}[Upper bound]\label{thm:upper}
$n[\alpha]=0$ in $H^2(\PN(T);\U(1))$.
\end{theorem}

\begin{proof}
The class $n[\alpha]$ is (property (C1), Baer sum) represented by the pushout of
\eqref{eq:ext1} along the $n$-th power map $[n]\colon\U(1)\to\U(1)$, $z\mapsto z^n$.
Concretely,
\[
E_n=\big(\U(1)\times\Nt\big)\big/\!\sim,\qquad
(z\lambda^{-n},\lambda g)\sim(z,g)\ \ (\lambda\in\U(1)),
\]
where $\U(1)\hookrightarrow\Nt$ as scalars; write $[z,g]$ for classes. The quotient map
$E_n\to\PN(T)$, $[z,g]\mapsto \bar g$, is a central extension with kernel
$\{[z,I]:z\in\U(1)\}\cong\U(1)$, and its class is $n[\alpha]$. Define
\[
\sigma\colon\PN(T)\longrightarrow E_n,\qquad \sigma(\bar g)=\big[(\det g)^{-1},\,g\big].
\]
\emph{Well defined:} if $g'=\lambda g$ ($\lambda\in\U(1)$) then, using
Lemma~\ref{lem:det} and the relation,
\[
\big[(\det g')^{-1},g'\big]=\big[(\lambda^n\det g)^{-1},\lambda g\big]
=\big[\lambda^{-n}(\det g)^{-1},\lambda g\big]=\big[(\det g)^{-1},g\big].
\]
\emph{Homomorphism:} the product in $E_n$ is $[z,g][w,h]=[zw,gh]$ (the $\U(1)$-factor is
central), so with $\det$ multiplicative
\[
\sigma(\bar g)\sigma(\bar h)=\big[(\det g)^{-1}(\det h)^{-1},gh\big]
=\big[(\det gh)^{-1},gh\big]=\sigma(\overline{gh}).
\]
Thus $\sigma$ is a continuous homomorphic splitting of $E_n\to\PN(T)$, so $n[\alpha]=0$.
\end{proof}
\STAT{\COMPLETE.} The section is genuinely a Lie-group homomorphism (no measurability
subtlety), being polynomial in matrix entries. This is the group-level shadow of the
$n$-torsion of Dixmier--Douady classes of $\PU(n)$-bundles \cite{DixmierDouady}, but the
proof above is self-contained.

\subsection{Lower bound: $k[\alpha]\neq0$ for $1\le k<n$}

The lower bound is proved by restricting to a \emph{finite} subgroup, where no continuity or
measurability subtlety can arise and all homomorphisms to $\U(1)$ are classical characters.

\begin{definition}[Finite scalar shadow]\label{def:shadow}
Let $\mu_n=\{z\in\U(1):z^n=1\}$, $D=\mu_n^{\,n}\subset T$ (diagonal matrices with entries in
$\mu_n$), and $Z=\{(\zeta,\dots,\zeta):\zeta\in\mu_n\}=D\cap\U(1)\cong\mu_n$ the scalar
subgroup. Put $W_n=D\rtimes S_n=\mu_n\wr S_n\subset\Nt$ and
$\Wn=W_n/Z=(D/Z)\rtimes S_n\subset\PN(T)$.
\end{definition}

\begin{lemma}[Restriction of {$[\alpha]$} to $\Wn$]\label{lem:restrict}
The restriction $[\alpha]|_{\Wn}\in H^2(\Wn;\U(1))$ is the image, under
$\mu_n\hookrightarrow\U(1)$, of the class of the finite central extension
\begin{equation}\label{eq:finshadow}
1\longrightarrow \mu_n\longrightarrow W_n\longrightarrow \Wn\longrightarrow 1 .
\end{equation}
\end{lemma}
\begin{proof}
The preimage of $\Wn$ under $\Nt\to\PN(T)$ is $\U(1)\cdot W_n$, and
$\U(1)\cdot W_n\cong(\U(1)\times W_n)/\{(\zeta^{-1},\zeta):\zeta\in\mu_n=\U(1)\cap W_n\}$ is
exactly the pushout of \eqref{eq:finshadow} along $\mu_n\hookrightarrow\U(1)$. Restriction
of extension classes is restriction of the group, giving the claim.
\end{proof}
\STAT{\COMPLETE.}

\begin{lemma}[$S_n$-equivariant homomorphism lemma]\label{lem:equiv}
Let $S_n$ act on $\mu_n^{\,n}$ (equivalently on $\Z_n^{\,n}$) by permuting coordinates and
trivially on $\U(1)$. Every $S_n$-invariant character
$\rho\colon\mu_n^{\,n}\to\U(1)$ has the form
\[
\rho\big(\zeta_n^{a_1},\dots,\zeta_n^{a_n}\big)=\zeta_n^{\,k\sum_i a_i},\qquad k\in\Z_n ,
\]
($\zeta_n=e^{2\pi i/n}$); there are exactly $n$ of them, and each satisfies
$\rho(\zeta_n,\dots,\zeta_n)=\zeta_n^{\,kn}=1$. Equivalently, the only $S_n$-equivariant
homomorphisms $\Z_n^{\,n}\to\Z_n$ are $x\mapsto k\sum_i x_i$, and all annihilate the
diagonal generator.
\end{lemma}
\begin{proof}
A character of the finite abelian group $\mu_n^{\,n}$ is $x\mapsto\prod_i\rho(e_i)^{a_i}$
with $\rho(e_i)\in\mu_n$; write $\rho(e_i)=\zeta_n^{c_i}$. Since $S_n$ permutes the standard
generators $e_i$ transitively, invariance $\rho(\sigma\cdot x)=\rho(x)$ forces
$c_{\sigma(i)}=c_i$ for all $\sigma$, hence $c_1=\dots=c_n=:k$. Then
$\rho(x)=\zeta_n^{k\sum a_i}$ and on the diagonal $\sum a_i=n$, so
$\rho=\zeta_n^{kn}=1$. There are exactly $n$ choices $k\in\Z_n$.
\end{proof}
\STAT{\COMPLETE.} This is the exact point at which ``$S_n$-equivariance'' enters: it forces
the invariant character to be a multiple of the coordinate sum, which is $\equiv0$ on the
diagonal modulo $n$.

\begin{theorem}[Lower bound]\label{thm:lower}
For every integer $k$ with $1\le k<n$, $k[\alpha]\neq0$ in $H^2(\PN(T);\U(1))$.
\end{theorem}

\begin{proof}
Restriction $H^2(\PN(T);\U(1))\to H^2(\Wn;\U(1))$ is a homomorphism, so it suffices to show
$k[\alpha]|_{\Wn}\neq0$ for $1\le k<n$. By Lemma~\ref{lem:restrict}, $k[\alpha]|_{\Wn}$ is
the class of the pushout of \eqref{eq:finshadow} along $\mu_n\xrightarrow{\ z\mapsto z^{k}\ }
\U(1)$. This pushout \emph{splits} (i.e.\ the class vanishes) if and only if there is a
homomorphism $\rho\colon W_n\to\U(1)$ extending $\zeta\mapsto\zeta^{k}$ on
$\mu_n=Z$, i.e.\ with
\begin{equation}\label{eq:need}
\rho(\zeta_n I)=\zeta_n^{\,k}=e^{2\pi i k/n}\ (\neq 1\text{ since }1\le k<n).
\end{equation}
Now for \emph{any} homomorphism $\rho\colon W_n=D\rtimes S_n\to\U(1)$, the restriction
$\rho|_{D}$ is $S_n$-invariant: for $s\in S_n\subset W_n$ and $d\in D$ we have
$s d s^{-1}=s\cdot d\in D$ and, $\U(1)$ being abelian,
$\rho(s\cdot d)=\rho(sds^{-1})=\rho(d)$. By Lemma~\ref{lem:equiv},
$\rho|_D(\zeta_n I)=\rho(\zeta_n,\dots,\zeta_n)=1$, contradicting \eqref{eq:need}. Hence no
such $\rho$ exists, so $k[\alpha]|_{\Wn}\neq0$, so $k[\alpha]\neq0$.
\end{proof}
\STAT{\COMPLETE.}

\begin{corollary}[Order exactly $n$]\label{thm:order}
The canonical class $[\alpha]$ --- equivalently $[\omega]\in H^2([X/\PU(n)];\underline{\U(1)})$
--- has order exactly $n$.
\end{corollary}
\begin{proof}
$n[\alpha]=0$ (Theorem~\ref{thm:upper}) and $k[\alpha]\neq0$ for $1\le k<n$
(Theorem~\ref{thm:lower}).
\end{proof}
\STAT{\COMPLETE.}

\subsection{Why no ``exotic'' splitting evades the argument}\label{ssec:exotic}

Three objections must be ruled out; all three are handled by the restriction to the
\emph{finite} group $\Wn$.

\begin{proposition}[Robustness of the lower bound]\label{prop:robust}
The conclusion $k[\alpha]\neq0$ ($1\le k<n$) is immune to:
\begin{enumerate}
\item[(a)] \textbf{Nonlinear ``splittings.''} A splitting of a central extension of groups
is by definition a group homomorphism $s$ with $\pi\circ s=\id$; ``nonlinear'' set-maps are
irrelevant, since vanishing of the class \emph{means} existence of such a homomorphism. The
proof of Theorem~\ref{thm:lower} shows no \emph{homomorphism} $\rho$ with the required value
exists.
\item[(b)] \textbf{Discontinuous homomorphisms.} Over the ambient Lie group $\PN(T)$ there
exist, by a Hamel-basis/choice construction, discontinuous homomorphisms $\PN(T)\to\U(1)$.
But any hypothetical splitting of $k[\alpha]$ over $\PN(T)$ restricts to a splitting over the
\emph{finite} subgroup $\Wn$, where every homomorphism is automatically continuous and is a
classical character. Theorem~\ref{thm:lower} rules those out. Formally: restriction
$H^2(\PN(T);\U(1))\to H^2(\Wn;\U(1))$ is a homomorphism, and $k[\alpha]\mapsto k[\alpha]
|_{\Wn}\neq0$, so $k[\alpha]\neq0$ regardless of the regularity of any putative section
upstairs.
\item[(c)] \textbf{Non-module (twisted) splittings.} The coefficient module is
\emph{central} ($\U(1)$ with trivial action); a section respecting only some twisted action
is not a splitting of \eqref{eq:ext1}. Since \eqref{eq:ext1} is central, the notion of
splitting is unambiguous and is the one used.
\end{enumerate}
\end{proposition}
\begin{proof}
Immediate from the definitions and from the naturality of restriction, as spelled out in
(a)--(c).
\end{proof}
\STAT{\COMPLETE.}

\begin{remark}[Where the upper and lower bounds meet]
The upper bound used a specific, everywhere-defined homomorphic section (the determinant
twist), so it is unconditional. The lower bound is a nonexistence statement pushed down to a
finite group, so it too is unconditional. Together they pin the order at exactly $n$ with no
appeal to any computer search; the repository's \texttt{thmG\_general.py} merely
double-checks Lemma~\ref{lem:equiv} and the order over $n=2,\dots,8$.
\end{remark}

% ======================================================================
\section{Item E: $H^2(\PN(T);\U(1))\cong\Z/n\oplus\Z_2$}\label{sec:E}
% ======================================================================

We now compute the full group, using the Lyndon--Hochschild--Serre (LHS) spectral sequence
of the split extension \eqref{eq:PNsplit}. All coefficients are $\U(1)$ with trivial action
(Definition~\ref{def:coh}); the outer group $S_n$ is finite, so $H^{\bullet}(S_n;-)$ is
ordinary group cohomology, while the inner torus contributes smooth (Segal--Mitchison)
cohomology.

\subsection{The spectral sequence and its coefficient inputs}

\begin{theorem}[Existence of the LHS spectral sequence]\label{thm:LHS}
For the split extension \eqref{eq:PNsplit} there is a first-quadrant cohomological spectral
sequence
\[
E_2^{p,q}=H^{p}\big(S_n;\,H^{q}_{\sm}(\Tp;\U(1))\big)\ \Longrightarrow\
H^{p+q}_{\sm}(\PN(T);\U(1)),
\]
with the usual multiplicative and edge structure, converging to a finite filtration of the
abutment.
\end{theorem}
\STAT{\SKETCH.} This is the Hochschild--Serre spectral sequence for measurable/smooth
cohomology of the extension of locally compact groups $\Tp\to\PN(T)\to S_n$, due to Moore
\cite[\S\S2--3]{MooreIV} (see also \cite{Lichtenbaum,Flach} for the topological/site
formulation); the inner terms $H^q_{\sm}(\Tp;\U(1))$ carry the $S_n$-action induced by
conjugation, and because $S_n$ is finite and discrete the outer cohomology is ordinary group
cohomology. We invoke Moore's theorem; its hypotheses (a closed normal subgroup of a second
countable locally compact group) hold.

We now identify each relevant coefficient group.

\begin{lemma}[Torus coefficients]\label{lem:toruscoh}
As $S_n$-modules,
\[
H^{0}_{\sm}(\Tp;\U(1))=\U(1)\ (\text{trivial}),\qquad
H^{1}_{\sm}(\Tp;\U(1))=\Thatp=A_{n-1},\qquad
H^{2}_{\sm}(\Tp;\U(1))=0 .
\]
\end{lemma}
\begin{proof}
Degree $0$ is the constants. Degree $1$: $H^1_{\sm}(\Tp;\U(1))=\Hom_{\mathrm{cont}}
(\Tp,\U(1))=\Thatp=A_{n-1}$ (Convention~\ref{conv:torus}), with the permutation
$S_n$-action.

Degree $2$: We claim $H^2_{\sm}(\Tp;\U(1))=0$, i.e.\ every smooth $\U(1)$-central extension
of the torus $\Tp$ splits. Use the exponential sequence of trivial $\Tp$-modules
$0\to\Z\to\R\xrightarrow{\exp(2\pi i\,\cdot)}\U(1)\to0$ and its long exact sequence in
$H^{\bullet}_{\sm}(\Tp;-)$. For the compact group $\Tp$, continuous cohomology with $\R$
coefficients vanishes in positive degree, $H^{k}_{\sm}(\Tp;\R)=0$ for $k\ge1$ (invariant
integration/averaging \cite[Ch.~IX]{BorelWallach}, \cite{MooreIII}), so the connecting maps
give $H^{k}_{\sm}(\Tp;\U(1))\cong H^{k+1}_{\sm}(\Tp;\Z)$ for $k\ge1$. Finally, for a compact
Lie group $G$, $H^{k}_{\sm}(G;\Z)\cong H^{k}(BG;\Z)$ (Segal \cite{Segal}, Wigner
\cite{Wigner}); for $G=\Tp\cong\U(1)^{n-1}$, $BG=(\C P^{\infty})^{n-1}$ has
$H^{\bullet}(BG;\Z)=\Z[x_1,\dots,x_{n-1}]$ with $\deg x_i=2$, concentrated in even degrees
and torsion-free. Hence $H^{3}(B\Tp;\Z)=0$, so
$H^{2}_{\sm}(\Tp;\U(1))\cong H^{3}_{\sm}(\Tp;\Z)\cong H^3(B\Tp;\Z)=0$.
\end{proof}
\STAT{\SKETCH.} The vanishing $H^2_{\sm}(\Tp;\U(1))=0$ is reduced to the two named results
$H^{k}_{\sm}(\mathrm{compact};\R)=0$ ($k\ge1$) and $H^{\bullet}_{\sm}(G;\Z)\cong
H^{\bullet}(BG;\Z)$ (Segal--Wigner), whose hypotheses (compactness, Lie) hold. Intuitively:
a continuous bicharacter $\Tp\times\Tp\to\U(1)$ is trivial (each slice is a character, hence
locally constant integer-valued), so the extension is abelian, and abelian smooth extensions
of a torus by $\U(1)$ split by the $BG$-computation.

\begin{lemma}[$S_n$-cohomology inputs]\label{lem:Sncoh}
With trivial $\U(1)$-coefficients and the permutation module $A_{n-1}$:
\begin{enumerate}
\item[(i)] $H^{0}(S_n;A_{n-1})=(A_{n-1})^{S_n}=0$.
\item[(ii)] $H^{1}(S_n;A_{n-1})\cong\Z/n$.
\item[(iii)] $H^{2}(S_n;\U(1))\cong M(S_n)$, the Schur multiplier: $\Z_2$ for $n\ge4$ and
$0$ for $n\le3$.
\end{enumerate}
\end{lemma}
\begin{proof}
(i) An $S_n$-invariant vector of $\Z^n$ is a scalar multiple of $(1,\dots,1)$; requiring
coordinate sum $0$ forces it to be $0$.

(ii) Use the short exact sequence of $S_n$-modules
$0\to A_{n-1}\to\Z^n\xrightarrow{\ \Sigma\ }\Z\to0$, where $\Z^n$ is the permutation module
(i.e.\ $\Z[S_n/S_{n-1}]$) and $\Z$ is trivial. Its long exact sequence reads
\[
H^0(S_n;\Z^n)\xrightarrow{\ \Sigma_*\ }H^0(S_n;\Z)\xrightarrow{\ \delta\ }
H^1(S_n;A_{n-1})\longrightarrow H^1(S_n;\Z^n).
\]
By Shapiro's lemma $H^{\bullet}(S_n;\Z^n)=H^{\bullet}(S_n;\Z[S_n/S_{n-1}])=
H^{\bullet}(S_{n-1};\Z)$, so $H^0(S_n;\Z^n)=\Z\cdot(1,\dots,1)$ and
$H^1(S_n;\Z^n)=H^1(S_{n-1};\Z)=\Hom(S_{n-1},\Z)=0$. The map $\Sigma_*$ sends $(1,\dots,1)$
to $n\in\Z=H^0(S_n;\Z)$, i.e.\ is multiplication by $n$; hence
$H^1(S_n;A_{n-1})\cong\mathrm{coker}(\Z\xrightarrow{n}\Z)=\Z/n$.

(iii) For finite $G$ and divisible $\U(1)$, the universal-coefficient/Hopf formula gives
$H^2(G;\U(1))\cong\Hom(H_2(G;\Z),\U(1))\cong H_2(G;\Z)=M(G)$ (the last isomorphism because
$H_2(G;\Z)$ is finite abelian, hence non-canonically self-Pontryagin-dual, and
$\Ext(H_1(G),\U(1))=0$ as $\U(1)$ is injective). Schur's computation of the multiplier of
the symmetric group is $M(S_n)=\Z_2$ for $n\ge4$ and $0$ for $n\in\{2,3\}$
\cite[Ch.~2]{Karpilovsky}, \cite{Schur}.
\end{proof}
\STAT{(i),(ii) \COMPLETE{}; (iii) \SKETCH{} (Schur's multiplier of $S_n$, cited).}

\subsection{The $E_2$ page in total degree $2$, and all differentials}

Assembling Lemmas~\ref{lem:toruscoh}--\ref{lem:Sncoh}, the anti-diagonal $p+q=2$ of the
$E_2$ page is
\[
E_2^{0,2}=H^0\big(S_n;H^2_{\sm}(\Tp;\U(1))\big)=H^0(S_n;0)=0,\qquad
E_2^{1,1}=H^1(S_n;A_{n-1})=\Z/n,
\]
\[
E_2^{2,0}=H^2(S_n;\U(1))=\begin{cases}\Z_2,& n\ge4,\\[1mm]0,& n\le3,\end{cases}
\]
and, in the neighbouring spots we shall need,
\[
E_2^{0,1}=H^0(S_n;A_{n-1})=0,\qquad E_2^{3,0}=H^3(S_n;\U(1)).
\]

\begin{proposition}[All differentials touching total degree $2$ vanish]\label{prop:diffs}
In the spectral sequence of Theorem~\ref{thm:LHS}:
\begin{enumerate}
\item[(i)] every differential \emph{into} the bottom row $E_r^{\bullet,0}$ vanishes; in
particular $E_\infty^{p,0}=E_2^{p,0}=H^p(S_n;\U(1))$ for all $p$;
\item[(ii)] consequently $d_2\colon E_2^{1,1}\to E_2^{3,0}$ is zero, and since no
differential enters $E^{1,1}$, we get $E_\infty^{1,1}=E_2^{1,1}=\Z/n$;
\item[(iii)] $E_\infty^{0,2}=E_2^{0,2}=0$ and $E_\infty^{2,0}=E_2^{2,0}$.
\end{enumerate}
\end{proposition}
\begin{proof}
(i) The extension \eqref{eq:PNsplit} is \emph{split}: the section
$j\colon S_n\hookrightarrow\PN(T)$ (permutation matrices) satisfies $p\circ j=\id_{S_n}$
where $p\colon\PN(T)\to S_n$. With trivial coefficients $\U(1)$, inflation
$\inf=p^{*}\colon H^{\bullet}(S_n;\U(1))\to H^{\bullet}(\PN(T);\U(1))$ and restriction
$\res=j^{*}$ satisfy $\res\circ\inf=(p\circ j)^{*}=\id$. Hence inflation is (split)
injective in every degree. The image of inflation is precisely the bottom filtration piece
$F^{p}H^{p}=E_\infty^{p,0}$, and the edge map $E_2^{p,0}\twoheadrightarrow E_\infty^{p,0}$
composed with $E_\infty^{p,0}\hookrightarrow H^p$ is inflation. Injectivity of inflation
forces the surjection $E_2^{p,0}\twoheadrightarrow E_\infty^{p,0}$ to be an isomorphism, so
no differential may hit the bottom row.

(ii) The only differential into $E^{3,0}$ that could reduce it is
$d_2\colon E_2^{1,1}\to E_2^{3,0}$ (and $d_3\colon E_3^{0,2}\to E_3^{3,0}$); by (i) both
vanish, in particular $d_2|_{E_2^{1,1}}=0$. No differential enters $E^{1,1}$ (sources
$E_r^{1-r,r}$ vanish for $r\ge2$ as $1-r<0$), and the only possible outgoing differential
$d_2\colon E_2^{1,1}\to E_2^{3,0}$ is zero; hence $E_\infty^{1,1}=E_2^{1,1}=\Z/n$.

(iii) $E_2^{0,2}=0$ has nothing in it; and the only differential into $E^{2,0}$ is
$d_2\colon E_2^{0,1}\to E_2^{2,0}$ with $E_2^{0,1}=0$, so $E_\infty^{2,0}=E_2^{2,0}$ (also a
special case of (i)).
\end{proof}
\STAT{\COMPLETE.} The vanishing of the crucial differential $d_2|_{E^{1,1}}$ is forced by
the \emph{splitting} of \eqref{eq:PNsplit} alone; it does \emph{not} require the order
theorem.

\subsection{The extension problem, and the answer}

\begin{theorem}[Full second cohomology]\label{thm:h2}
For all $n\ge2$,
\[
H^2(\PN(T);\U(1))\ \cong\ \Z/n\ \oplus\ H^2(S_n;\U(1))\ =\
\begin{cases}
\Z/n\oplus\Z_2, & n\ge4,\\
\Z/3, & n=3,\\
\Z/2, & n=2,
\end{cases}
\]
and the canonical class $[\alpha]$ maps to a generator of the $\Z/n$ summand.
\end{theorem}
\begin{proof}
By Proposition~\ref{prop:diffs} the degree-$2$ filtration of $H^2:=H^2(\PN(T);\U(1))$ is
\[
0=F^{3}\subseteq F^{2}\subseteq F^{1}=F^{0}=H^2,\qquad
F^{2}=E_\infty^{2,0}=H^2(S_n;\U(1)),\quad F^1/F^2=E_\infty^{1,1}=\Z/n,\quad
E_\infty^{0,2}=0,
\]
i.e.\ a short exact sequence
\begin{equation}\label{eq:extprob}
0\longrightarrow H^2(S_n;\U(1))\xrightarrow{\ \inf\ }H^2\xrightarrow{\ \rho\ }\Z/n
\longrightarrow0 .
\end{equation}
\emph{The sequence splits, for every $n$.} The sub $F^2=\im(\inf)$ is the image of
inflation, and by Proposition~\ref{prop:diffs}(i) the restriction
$\res\colon H^2\to H^2(S_n;\U(1))$ satisfies $\res\circ\inf=\id$. Thus $\res$ is a retraction
of the inclusion $\inf$ in \eqref{eq:extprob}, so
\[
H^2=\im(\inf)\oplus\ker(\res)\cong H^2(S_n;\U(1))\oplus\big(H^2/F^2\big)
=H^2(S_n;\U(1))\oplus\Z/n .
\]
This gives the stated groups using Lemma~\ref{lem:Sncoh}(iii) and the low-$n$ values.

\emph{$[\alpha]$ generates the $\Z/n$.} By Corollary~\ref{thm:order}, $[\alpha]$ has order
exactly $n$. Its image $\rho([\alpha])\in\Z/n$ has order $n/|\langle[\alpha]\rangle\cap F^2|$;
since $F^2=H^2(S_n;\U(1))$ has order dividing $2$, and $\ker(\res)$ is a complement with
$\res([\alpha])\in H^2(S_n;\U(1))$, we argue directly: the projection $q:=$
($H^2\twoheadrightarrow\ker(\res)\cong\Z/n$) along $F^2$ agrees with $\rho$ up to the chosen
splitting, and $q([\alpha])$ has order equal to that of $[\alpha]$ minus its $F^2$-part.
Concretely, write $[\alpha]=(u,v)\in F^2\oplus\Z/n$. Then $\res([\alpha])=u$, and by
Proposition~\ref{prop:robust} applied through the finite shadow, $k[\alpha]\neq0$ for
$1\le k<n$; if $v$ had order $d<n$ in $\Z/n$ then $d\,[\alpha]=(du,0)=(du,0)$ would lie in
$F^2$, of order $\le2$, forcing $2d\,[\alpha]=0$, i.e.\ $n\mid 2d$. Combined with $d\mid n$
and $d<n$ this leaves only $n$ even and $d=n/2$. That residual case is excluded in
Proposition~\ref{prop:gen} below; granting it, $v$ has order $n$, i.e.\ generates.
\end{proof}
\STAT{\COMPLETE{} for the group isomorphism $H^2\cong\Z/n\oplus H^2(S_n;\U(1))$ (the
splitting of \eqref{eq:extprob} is unconditional, via the retraction $\res$); the
identification of $[\alpha]$ as a \emph{generator} of the $\Z/n$ summand is \COMPLETE{} for
$n$ odd and rests on Proposition~\ref{prop:gen} for $n$ even.}

\begin{proposition}[{$[\alpha]$} generates modulo the Schur summand]\label{prop:gen}
$\rho([\alpha])$ generates $E_\infty^{1,1}=\Z/n$.
\begin{itemize}
\item[$\bullet$] For $n$ odd this is \COMPLETE: $|\langle[\alpha]\rangle|=n$ is coprime to
$|F^2|\in\{1,2\}$, so $\langle[\alpha]\rangle\cap F^2=0$ and $\rho([\alpha])$ has order $n$.
\item[$\bullet$] For $n$ even this is a \GAP: one must exclude $(n/2)[\alpha]\in F^2$, i.e.\
that the order-$2$ element $(n/2)[\alpha]$ is inflated from the Schur class of $S_n$.
\end{itemize}
\end{proposition}
\STAT{$n$ odd: \COMPLETE. $n$ even: \GAP{} (stated precisely below).}

\begin{remark}[Exact statement of the residual gap, and partial evidence]\label{rem:gap}
The residual open sub-lemma is:
\begin{quote}\itshape
For $n$ even, $(n/2)[\alpha]\notin\im\big(\inf\colon H^2(S_n;\U(1))\to
H^2(\PN(T);\U(1))\big)$; equivalently, $\res_{S_n}\big((n/2)[\alpha]\big)=0$ does \emph{not}
coincide with $(n/2)[\alpha]$ being the inflated Schur class.
\end{quote}
Two remarks bound the stakes. (1) The isomorphism type $H^2\cong\Z/n\oplus\Z_2$ is
\emph{already proved unconditionally} (Theorem~\ref{thm:h2}); the gap concerns only whether
the \emph{specific} class $[\alpha]$ is a generator of the $\Z/n$ factor versus lying in
$\Z/n\oplus\Z_2$ with a nonzero $\Z_2$-component (which would still leave $[\alpha]$ of order
$n$ and the group unchanged). (2) The natural attempt --- detect a nonzero image on the
abelian part $D/Z\subset\Wn$ --- fails because $[\alpha]|_{D/Z}=0$: the finite abelian
extension $1\to Z\to D\to D/Z\to1$ splits (the retraction $D=\mu_n^n\to\mu_n$,
$x\mapsto x_1$, restricts to the identity on $Z$), so $[\alpha]$ is invisible on $D/Z$ and
lives genuinely in the mixed $E^{1,1}$-stratum. Closing the gap therefore requires computing
the $\U(1)$-valued $2$-cocycle of $\Wn$ on a permutation$\,\times\,$sign pair (the
transgression $S_n\to\Thatp$), or equivalently identifying $\res_{S_n}(n/2)[\alpha]$ with a
$2$-cocycle of $S_n$ and checking it is a coboundary. A direct computation of the
$E^{1,1}$-transgression cocycle $\gamma(\sigma)=e_{\sigma^{-1}(n)}-e_n\in A_{n-1}$ for the
section of Convention~\ref{conv:torus} strongly indicates $\gamma$ is a \emph{generator} of
$H^1(S_n;A_{n-1})=\Z/n$ (it is the connecting image of $1\in\Z$ under
$\Sigma\colon\Z^n\to\Z$), which would force $\rho([\alpha])$ to generate and close the gap;
we do not certify the cocycle-identity bookkeeping here and so leave it flagged.
\end{remark}

\subsection{Coefficient conventions, topology assumptions, low-$n$ exceptions}

\begin{itemize}
\item \textbf{Coefficients.} Everywhere $\U(1)$ is a \emph{trivial} module (central
extensions); $A_{n-1}=\Thatp$ carries the permutation $S_n$-action; $\Z$ and $\R$ in
Lemma~\ref{lem:toruscoh} are trivial modules.
\item \textbf{Topology.} $\Tp,\PN(T),\PU(n),\U(n)$ are compact Lie groups; ``$H^2_{\sm}$''
is Segal--Mitchison smooth cohomology, which for these (second countable, compact) groups
coincides with Moore measurable cohomology \cite{MooreIII,WagnerAustin}. The inner
cohomology $H^{\bullet}_{\sm}(\Tp;\U(1))$ is the smooth one; the outer $H^{\bullet}(S_n;-)$
is ordinary (discrete) group cohomology because $S_n$ is finite. The identification
$H^{\bullet}_{\sm}(\Tp;\Z)\cong H^{\bullet}(B\Tp;\Z)$ requires compact Lie $\Tp$
(Segal--Wigner).
\item \textbf{Low-$n$ exceptions.} $n=2$: $S_2=\Z_2$, $A_1=\Z$ with the sign action,
$H^1(\Z_2;\Z_{\mathrm{sgn}})=\Z/2$ and $H^2(\Z_2;\U(1))=0$, so $H^2(\PN(T);\U(1))=\Z/2$
(no Schur factor). $n=3$: $M(S_3)=0$, so $H^2(\PN(T);\U(1))=\Z/3$. The $\Z_2$ summand
appears only for $n\ge4$, exactly the range in which the Schur multiplier is nontrivial and
(Paper~A) the Peres--Mermin shadow is available.
\end{itemize}
\STAT{\COMPLETE{} (bookkeeping of conventions and the low-$n$ computations).}

% ======================================================================
\section{Status of each step}\label{sec:status}
% ======================================================================

The table records, for every lettered item and every numbered result, whether the argument
here is a \emph{complete proof} (checkable without repository scripts), a \emph{proof
sketch} reduced to a named standard theorem (with hypotheses verified), or an explicit
\emph{gap}.

\begin{center}
\small
\begin{tabular}{@{}p{0.10\textwidth}p{0.53\textwidth}p{0.30\textwidth}@{}}
\toprule
\textbf{Item / Result} & \textbf{Content} & \textbf{Status} \\
\midrule
A & Objects; $X\cong\PU(n)/\PN(T)$; choice of cohomology theory & Complete (defs);
Sketch for (C1)--(C3) [BX, TXL, Seg, Moore] \\
\;Lem.~\ref{lem:homog} & Structure/homogeneity of $X$ & Complete \\
\;Def.~\ref{def:coh} & Groupoid cohomology $H^\bullet(-;\underline{\U(1)})$ & Complete def.;
Sketch (C1)--(C3) \\
\midrule
B & Canonical class; vacuity; nontriviality & Complete (given C, D) \\
\;Thm.~\ref{thm:vacuity} & Vacuity on $\U(n)\ltimes X$ & Complete \\
\;Thm.~\ref{thm:nontrivial} & Nontriviality on $\PU(n)\ltimes X$ & Complete given C, D \\
\midrule
D & Morita reduction to $\PN(T)$ & Complete bibundle; Sketch (C3) \\
\;Thm.~\ref{thm:MoritaGeneral} & $G\ltimes(G/H)\simeq_{\mathrm{Mor}}H$ & Sketch [MM, dH,
BX] \\
\;Prop.~\ref{prop:morita} & Bibundle $\PU(n)$; class $\mapsto$ \eqref{eq:ext1} & Complete
(constr.\ \& transport); Sketch (C3) \\
\midrule
C & Order exactly $n$ & Complete \\
\;Thm.~\ref{thm:upper} & $n[\alpha]=0$ (determinant twist) & Complete \\
\;Lem.~\ref{lem:equiv} & $S_n$-equivariant hom lemma & Complete \\
\;Thm.~\ref{thm:lower} & $k[\alpha]\neq0$, $1\le k<n$ & Complete \\
\;Cor.~\ref{thm:order} & Order $=n$ & Complete \\
\;Prop.~\ref{prop:robust} & No nonlinear/discontinuous/non-module evasion & Complete \\
\midrule
E & $H^2(\PN(T);\U(1))\cong\Z/n\oplus\Z_2$ ($n\ge4$) & Complete (iso); Sketch inputs;
Gap (generator, $n$ even) \\
\;Thm.~\ref{thm:LHS} & LHS spectral sequence exists & Sketch [Moore IV] \\
\;Lem.~\ref{lem:toruscoh} & $H^{\le2}_{\sm}(\Tp;\U(1))$ & $q{=}0,1$ Complete; $q{=}2{=}0$
Sketch [Seg, Wig] \\
\;Lem.~\ref{lem:Sncoh} & $H^0{=}0$, $H^1{=}\Z/n$, $H^2(S_n;\U(1))$ & (i),(ii) Complete;
(iii) Sketch [Schur] \\
\;Prop.~\ref{prop:diffs} & All differentials to deg.\ $2$ vanish & Complete \\
\;Thm.~\ref{thm:h2} & $H^2\cong\Z/n\oplus H^2(S_n;\U(1))$ & Complete (group iso) \\
\;Prop.~\ref{prop:gen} & $[\alpha]$ generates $\Z/n$ & Complete ($n$ odd); \textbf{Gap}
($n$ even) \\
\bottomrule
\end{tabular}
\end{center}

\medskip
\noindent\textbf{Single most important gap to close first.}
Proposition~\ref{prop:gen}/Remark~\ref{rem:gap} for \emph{even} $n$: prove that
$(n/2)[\alpha]$ is not the inflated Schur class, equivalently that the transgression cocycle
$\gamma\in H^1(S_n;A_{n-1})=\Z/n$ of the extension \eqref{eq:ext1} is a generator. This does
not affect the isomorphism type of $H^2$ (already proved), only the statement that the
\emph{canonical} class is a generator of the $\Z/n$ summand rather than a generator-plus-Schur.

% ======================================================================
\begin{thebibliography}{99}

\bibitem{BehrendXu} K.~Behrend, P.~Xu, \emph{Differentiable stacks and gerbes},
J.~Symplectic Geom.\ \textbf{9} (2011), 285--341.

\bibitem{TXL} J.-L.~Tu, P.~Xu, C.~Laurent-Gengoux, \emph{Twisted $K$-theory of
differentiable stacks}, Ann.\ Sci.\ \'Ecole Norm.\ Sup.\ (4) \textbf{37} (2004), 841--910.

\bibitem{Crainic} M.~Crainic, \emph{Differentiable and algebroid cohomology, van Est
isomorphisms, and characteristic classes}, Comment.\ Math.\ Helv.\ \textbf{78} (2003),
681--721.

\bibitem{Segal} G.~Segal, \emph{Cohomology of topological groups}, in \emph{Symposia
Mathematica IV} (INDAM, 1968/69), Academic Press (1970), 377--387.

\bibitem{MooreIII} C.~C.~Moore, \emph{Group extensions and cohomology for locally compact
groups. III}, Trans.\ Amer.\ Math.\ Soc.\ \textbf{221} (1976), 1--33.

\bibitem{MooreIV} C.~C.~Moore, \emph{Group extensions and cohomology for locally compact
groups. IV}, Trans.\ Amer.\ Math.\ Soc.\ \textbf{221} (1976), 35--58.

\bibitem{Wigner} D.~Wigner, \emph{Algebraic cohomology of topological groups}, Trans.\
Amer.\ Math.\ Soc.\ \textbf{178} (1973), 83--93.

\bibitem{WagnerAustin} T.~Austin, C.~C.~Moore, \emph{Continuity properties of measurable
group cohomology}, Math.\ Ann.\ \textbf{356} (2013), 885--937.

\bibitem{BorelWallach} A.~Borel, N.~Wallach, \emph{Continuous Cohomology, Discrete
Subgroups, and Representations of Reductive Groups}, 2nd ed., Math.\ Surveys Monogr.\
\textbf{67}, Amer.\ Math.\ Soc.\ (2000).

\bibitem{Lichtenbaum} S.~Lichtenbaum, \emph{The Weil-\'etale topology on schemes over finite
fields}, Compos.\ Math.\ \textbf{141} (2005), 689--702. (Hochschild--Serre for site
cohomology.)

\bibitem{Flach} M.~Flach, \emph{Cohomology of topological groups with applications to the
Weil group}, Compos.\ Math.\ \textbf{144} (2008), 633--656.

\bibitem{MoerdijkMrcun} I.~Moerdijk, J.~Mr\v{c}un, \emph{Introduction to Foliations and Lie
Groupoids}, Cambridge Stud.\ Adv.\ Math.\ \textbf{91}, Cambridge Univ.\ Press (2003).

\bibitem{delHoyo} M.~del Hoyo, \emph{Lie groupoids and their orbispaces}, Port.\ Math.\
\textbf{70} (2013), 161--209.

\bibitem{Karpilovsky} G.~Karpilovsky, \emph{The Schur Multiplier}, London Math.\ Soc.\
Monogr.\ (N.S.)\ \textbf{2}, Oxford Univ.\ Press (1987).

\bibitem{Schur} I.~Schur, \emph{\"Uber die Darstellung der symmetrischen und der
alternierenden Gruppe durch gebrochene lineare Substitutionen}, J.~Reine Angew.\ Math.\
\textbf{139} (1911), 155--250.

\bibitem{DixmierDouady} J.~Dixmier, A.~Douady, \emph{Champs continus d'espaces hilbertiens
et de $C^*$-alg\`ebres}, Bull.\ Soc.\ Math.\ France \textbf{91} (1963), 227--284.

\bibitem{Brown} K.~S.~Brown, \emph{Cohomology of Groups}, Grad.\ Texts in Math.\
\textbf{87}, Springer (1982). (LHS spectral sequence, Shapiro's lemma, inflation--restriction.)

\end{thebibliography}

\end{document}
